% Candidate standalone chapter for GAGA source units gaga:no:1 through
% gaga:no:4. This is deliberately the reduced 1956 notion. It does not
% introduce nonreduced complex analytic spaces.

\input{preamble}

% OK, start here.
%
\begin{document}

\title{Complex Analytic Spaces and GAGA}

\maketitle

\phantomsection
\label{section-phantom}

\tableofcontents

\section{Introduction}
\label{section-introduction}

\noindent
This chapter develops the complex-analytic foundations needed to state and
prove the comparison theorems usually called GAGA. The initial sections use
the reduced notion of complex analytic space in
\cite[\S 1, nos. 1--4, pp. 3--7]{GAGA}. This scope is recorded explicitly
because a structure sheaf contained in a sheaf of functions has no
nilpotents. The abstract language of ringed spaces and coherent modules is
that of Sheaves, Definition \ref{sheaves-definition-ringed-space} and
Modules, Section \ref{modules-section-coherent}.

\section{Analytic subsets of complex affine space}
\label{section-analytic-subsets}

\begin{definition}
\label{definition-analytic-subset}
\begin{reference}
\cite[\S 1, no. 1, pp. 3--4]{GAGA}
\end{reference}
Let $n \geq 0$ and endow $\mathbf C^n$ with its usual topology. A subset
$U \subset \mathbf C^n$ is an {\it analytic subset in the sense of Serre}
if, for every $x \in U$, there are an open neighbourhood $W$ of $x$ and
holomorphic functions $f_1, \ldots, f_k$ on $W$ such that
$$
U \cap W = \{z \in W \mid f_1(z)=\ldots=f_k(z)=0\}.
$$

Let $\mathcal C_U$ be the sheaf of germs of all complex-valued functions on
$U$. Let $\mathcal O_{\mathbf C^n}$ be the sheaf of holomorphic functions on
$\mathbf C^n$. The {\it sheaf of holomorphic functions on $U$} is the image
sheaf
$$
\mathcal O_U =
\mathop{\rm Im}\left(
\mathcal O_{\mathbf C^n}|_U \longrightarrow \mathcal C_U
\right).
$$
Thus, if $x \in U$ and $\mathcal I_{U,x}$ is the ideal of germs in
$\mathcal O_{\mathbf C^n,x}$ which vanish on $U$ near $x$, then
$$
\mathcal O_{U,x}=
\mathcal O_{\mathbf C^n,x}/\mathcal I_{U,x}.
$$
\end{definition}

\begin{definition}
\label{definition-holomorphic-map-analytic-subsets}
\begin{reference}
\cite[\S 1, no. 1, pp. 3--4]{GAGA}
\end{reference}
Let $U \subset \mathbf C^r$ and $V \subset \mathbf C^s$ be analytic
subsets. A map $\varphi:U\to V$ is {\it holomorphic} if it is continuous and
composition with $\varphi$ sends every germ in
$\mathcal O_{V,\varphi(x)}$ to a germ in $\mathcal O_{U,x}$ for every
$x\in U$. Equivalently, the $s$ coordinate functions of $\varphi$ are
sections of $\mathcal O_U$ locally on $U$.

A bijective holomorphic map is an {\it analytic isomorphism} if its inverse
is holomorphic.
\end{definition}

\begin{lemma}
\label{lemma-analytic-subset-basic-properties}
Let $U \subset \mathbf C^n$ and $U' \subset \mathbf C^{n'}$ be analytic
subsets in the sense of Definition \ref{definition-analytic-subset}.
\begin{enumerate}
\item The subset $U$ is locally closed and locally compact.
\item The stalk $\mathcal O_{U,x}$ is a reduced local $\mathbf C$-algebra
whose residue homomorphism is evaluation at $x$.
\item The subset $U\times U'$ is analytic in $\mathbf C^{n+n'}$, its
topology is the product topology, and its structure sheaf is the one obtained
from the product analytic charts.
\item Products of holomorphic maps are holomorphic.
\end{enumerate}
\end{lemma}

\begin{proof}
After shrinking around $x\in U$, the set $U$ is the common zero locus of
finitely many continuous functions and hence is closed in an open subset of
$\mathbf C^n$. This proves (1). The displayed quotient in Definition
\ref{definition-analytic-subset} embeds into $\mathcal C_{U,x}$, and a germ
is invertible exactly when its value at $x$ is nonzero. This proves (2).

Near $(x,x')$, concatenate local defining equations for $U$ and $U'$.
This proves analyticity of the product. The remaining assertions in (3) and
(4) follow directly from the induced topology, the coordinate criterion in
Definition \ref{definition-holomorphic-map-analytic-subsets}, and composition
of holomorphic functions.
\end{proof}


\section{Serre complex analytic spaces}
\label{section-serre-analytic-spaces}

\begin{definition}
\label{definition-serre-analytic-space}
\begin{reference}
\cite[\S 1, no. 2, pp. 4--5]{GAGA}
\end{reference}
A {\it complex analytic space in the sense of Serre} is a Hausdorff
topological space $X$ together with a subsheaf of $\mathbf C$-algebras
$$
\mathcal O_X \subset \mathcal C_X
$$
such that $X$ has an open covering $X=\bigcup V_i$ for which every ringed
space $(V_i,\mathcal O_X|_{V_i})$ is analytically isomorphic to an analytic
subset of some $\mathbf C^{n_i}$ endowed with the sheaf of Definition
\ref{definition-analytic-subset}.

A map of such spaces is {\it holomorphic} if it is continuous and pullback
by composition sends germs of holomorphic functions to germs of holomorphic
functions. An isomorphism is a holomorphic map with holomorphic inverse.
\end{definition}

\begin{remark}[Reduced scope]
\label{remark-serre-analytic-spaces-reduced}
Because $\mathcal O_X$ is a subsheaf of a sheaf of functions, every
$\mathcal O_{X,x}$ is reduced. Thus Definition
\ref{definition-serre-analytic-space} is the reduced 1956 notion used by
Serre. It does not include nilpotent analytic thickenings. Any later extension
to nonreduced complex analytic spaces must be introduced separately and must
not be attributed to GAGA, Definition 1.
\end{remark}

\begin{definition}
\label{definition-analytic-subspace}
Let $X$ be a Serre complex analytic space. A subset $Y\subset X$ is an
{\it analytic subspace} if, in every analytic chart $V\to U$, the image of
$Y\cap V$ is an analytic subset of $U$. It has the induced analytic-space
structure. When $Y$ is closed in $X$, we call it a {\it closed analytic
subspace}.
\end{definition}

\begin{lemma}
\label{lemma-serre-analytic-subspaces-products}
\begin{reference}
\cite[\S 1, no. 2, p. 5]{GAGA}
\end{reference}
Let $X$ and $X'$ be Serre complex analytic spaces.
\begin{enumerate}
\item Every analytic subspace of $X$ is locally closed and the induced
structure is independent of the chosen charts.
\item There is a unique Serre complex analytic-space structure on
$X\times X'$ for which products of charts are charts. Its topology is the
product topology.
\item Serre complex analytic spaces and holomorphic maps form a category.
\end{enumerate}
\end{lemma}

\begin{proof}
All three assertions are local on the source spaces. They therefore follow
from Lemma \ref{lemma-analytic-subset-basic-properties}, compatibility of the
restriction maps defining the structure sheaves, and the coordinate
criterion for holomorphic maps.
\end{proof}


\section{Analytic modules and Oka--Cartan coherence}
\label{section-analytic-modules}

\begin{definition}
\label{definition-analytic-module}
\begin{reference}
\cite[\S 1, no. 3, p. 5]{GAGA}
\end{reference}
Let $X$ be a Serre complex analytic space. An {\it analytic module} on $X$
is an $\mathcal O_X$-module. It is {\it coherent} if it is coherent in the
sense of Modules, Definition \ref{modules-definition-coherent}.
\end{definition}

\begin{definition}
\label{definition-vanishing-ideal-analytic-subspace}
Let $i:Y\to X$ be a closed analytic subspace. The {\it vanishing ideal}
$\mathcal J_Y\subset\mathcal O_X$ has stalk
$$
(\mathcal J_Y)_x=
\{f\in\mathcal O_{X,x}\mid f|_Y=0\text{ near }x\}.
$$
There is a canonical isomorphism
$$
\mathcal O_X/\mathcal J_Y \longrightarrow i_*\mathcal O_Y.
$$
\end{definition}

\begin{theorem}[Oka--Cartan coherence]
\label{theorem-oka-cartan-coherence}
\begin{reference}
\cite[\S 1, Proposition 1, pp. 5--6]{GAGA}
\end{reference}
Let $X$ be a Serre complex analytic space.
\begin{enumerate}
\item The $\mathcal O_X$-module $\mathcal O_X$ is coherent.
\item If $Y\subset X$ is a closed analytic subspace, then its vanishing
ideal $\mathcal J_Y$ is a coherent $\mathcal O_X$-module.
\end{enumerate}
\end{theorem}

\begin{proof}
The analytic input is the Oka--Cartan theorem that, for an open subset
$\Omega\subset\mathbf C^n$, the sheaf $\mathcal O_\Omega$ and the vanishing
ideal of every closed analytic subset of $\Omega$ are coherent.

The assertions are local on $X$. Choose a chart and shrink it so that
$X=Z(\mathcal I_X)$ is a closed analytic subset of an open
$\Omega\subset\mathbf C^n$. If $Y\subset X$ is closed analytic, shrink again
and write $Y=Z(\mathcal I_Y)$ with
$\mathcal I_X\subset\mathcal I_Y\subset\mathcal O_\Omega$. Then
$$
\mathcal O_X=\mathcal O_\Omega/\mathcal I_X,
\qquad
\mathcal J_Y=\mathcal I_Y/\mathcal I_X.
$$
The Oka--Cartan input makes $\mathcal O_\Omega$, $\mathcal I_X$, and
$\mathcal I_Y$ coherent. Modules, Lemmas
\ref{modules-lemma-coherent-abelian} and
\ref{modules-lemma-coherent-change-rings} now give (1) and (2).
\end{proof}

\begin{remark}[Examples]
\label{remark-examples-coherent-analytic-modules}
The sheaf of holomorphic sections of a holomorphic vector bundle is finite
locally free and hence coherent by Theorem
\ref{theorem-oka-cartan-coherence}; compare Cohomology, Lemma
\ref{cohomology-lemma-holomorphic-vector-bundles-locally-free-analytic-modules}.
Serre also records the sheaves of germs of automorphic functions from the
Cartan seminar as coherent examples; this cited analytic construction is not
needed for the foundation statements above.
\end{remark}


\section{Local analytic algebras}
\label{section-local-analytic-algebras}

\begin{lemma}
\label{lemma-local-analytic-algebra}
\begin{reference}
\cite[\S 1, no. 4, p. 6]{GAGA}
\end{reference}
Let $X$ be a Serre complex analytic space and $x\in X$. There are an
integer $n\geq 0$ and a radical ideal
$\mathfrak a\subset\mathbf C\{z_1,\ldots,z_n\}$ such that
$$
\mathcal O_{X,x}\cong
\mathbf C\{z_1,\ldots,z_n\}/\mathfrak a
$$
as local $\mathbf C$-algebras. Its maximal ideal consists of germs vanishing
at $x$, its residue field is $\mathbf C$, and it is reduced and Noetherian.
\end{lemma}

\begin{proof}
Choose an analytic chart at $x$ and translate its image point to the origin.
Definition \ref{definition-analytic-subset} gives the displayed quotient.
The ideal is radical because it is the ideal of germs vanishing on a set.
The residue and reducedness statements also follow from Lemma
\ref{lemma-analytic-subset-basic-properties}. The remaining analytic input is
the Weierstrass theorem that
$\mathbf C\{z_1,\ldots,z_n\}$ is Noetherian; the quotient is then Noetherian
by Algebra, Lemma \ref{algebra-lemma-Noetherian-permanence}.
\end{proof}

\begin{lemma}
\label{lemma-analytic-germ-local-algebra}
\begin{reference}
\cite[\S 1, no. 4, p. 6]{GAGA}
\end{reference}
The local analytic $\mathbf C$-algebra $\mathcal O_{X,x}$ determines the
germ of a Serre complex analytic space $X$ at $x$. More precisely, an
isomorphism of local analytic $\mathbf C$-algebras
$\mathcal O_{X,x}\cong\mathcal O_{Y,y}$ is induced contravariantly by an
isomorphism of analytic germs $(X,x)\cong(Y,y)$.
\end{lemma}

\begin{proof}
Present both local algebras as quotients of convergent power-series rings.
Images of the coordinate germs under a local homomorphism are convergent
germs and hence define a holomorphic map germ. The relations force its image
to lie in the required analytic germ. Applying the same construction to the
inverse algebra homomorphism gives inverse holomorphic germs.
\end{proof}

\begin{lemma}
\label{lemma-simple-point-regular-local}
\begin{reference}
\cite[\S 1, no. 4, p. 6]{GAGA}
\end{reference}
Let $X$ be a Serre complex analytic space and $x\in X$. The following are
equivalent for $n\geq 0$:
\begin{enumerate}
\item the germ $(X,x)$ is isomorphic to $(\mathbf C^n,0)$;
\item $\mathcal O_{X,x}\cong\mathbf C\{z_1,\ldots,z_n\}$;
\item $\mathcal O_{X,x}$ is a regular local ring of dimension $n$ in the
sense of Algebra, Definition \ref{algebra-definition-regular-local}.
\end{enumerate}
Such an $x$ is a {\it simple point of dimension $n$}. If every point of $X$
is simple, then $X$ is an {\it analytic variety}.
\end{lemma}

\begin{proof}
The equivalence of (1) and (2) follows from Lemma
\ref{lemma-analytic-germ-local-algebra}. The implication (2) $\Rightarrow$
(3) is the standard dimension computation for convergent power-series
rings. The converse is the regular analytic local-algebra theorem, proved
from the analytic inverse-function theorem (equivalently, by choosing a
basis of $\mathfrak m_x/\mathfrak m_x^2$ and applying the analytic implicit
function theorem). This analytic theorem is an explicit external dependency
of the present foundation block.
\end{proof}

\begin{lemma}
\label{lemma-components-dimension-analytic-space}
\begin{reference}
\cite[\S 1, no. 4, pp. 6--7]{GAGA}
\end{reference}
Let $X$ be a Serre complex analytic space and $x\in X$.
\begin{enumerate}
\item The irreducible components $X_i$ of the germ $(X,x)$ correspond to
the finitely many minimal primes $\mathfrak p_i$ of $\mathcal O_{X,x}$, and
$$
0=\bigcap_i\mathfrak p_i,
\qquad
\mathcal O_{X_i,x}=\mathcal O_{X,x}/\mathfrak p_i.
$$
\item The analytic dimension of $X$ at $x$ (one half of its local topological
dimension) is the supremum of the dimensions of the $X_i$ and equals the
Krull dimension of $\mathcal O_{X,x}$.
\end{enumerate}
\end{lemma}

\begin{proof}
Reducedness and Noetherianity are given by Lemma
\ref{lemma-local-analytic-algebra}. The algebraic assertion about the
finiteness and intersection of minimal primes follows from Algebra, Lemmas
\ref{algebra-lemma-Noetherian-irreducible-components} and
\ref{algebra-lemma-Zariski-topology}. The correspondence with analytic
components uses the local analytic Nullstellensatz. For (2), the analytic
normalization/parameter theorem says that an irreducible analytic local
algebra of analytic dimension $r$ is finite over
$\mathbf C\{z_1,\ldots,z_r\}$. Dimension invariance under finite integral
extensions then gives Krull dimension $r$. The analytic Nullstellensatz and
local parameter theorem are explicit external dependencies; no algebraic
Stacks result alone supplies them.
\end{proof}




\input{chapters}

\bibliography{my}
\bibliographystyle{amsalpha}

\end{document}
